\subsection*{Appendix E: The Reserve and Runway Policy Memo}

\begin{quote}
\small
\textbf{ENTITY AND BANK ACCOUNTS INCLUDED:}\\
\textbf{WHICH CASH IS RESTRICTED OR NOT ACCESSIBLE:}\\
\textbf{WHICH CUSTOMER OBLIGATIONS REMAIN UNFULFILLED:}\\
\textbf{UNAVOIDABLE PAYROLL AND SUPPLIER PAYMENTS:}\\
\textbf{DEBT SERVICE AND COVENANT SCHEDULE:}\\
\textbf{TAX AND STATUTORY OBLIGATIONS:}\\
\textbf{LIKELY REFUND / REWORK / INTERRUPTION EXPOSURE:}\\
\textbf{WORST CREDIBLE COLLECTION DELAY:}\\
\textbf{TIME REQUIRED TO REDUCE COSTS RESPONSIBLY:}\\
\textbf{BASE-CASE MINIMUM WEEKLY CASH:}\\
\textbf{ADVERSE-CASE MINIMUM WEEKLY CASH:}\\
\textbf{SEVERE-CASE MINIMUM WEEKLY CASH:}\\
\textbf{MINIMUM RESERVE FLOOR, WITH RATIONALE:}\\
\textbf{EARLY WARNING TRIGGER:}\\
\textbf{HARD STOP ON NEW DISCRETIONARY COMMITMENTS:}\\
\textbf{PERSON AUTHORIZED TO OVERRIDE, AND DOCUMENTATION:}\\
\textbf{NEXT REVIEW DATE:}
\end{quote}

Do not automatically copy a 3- or 6-month reserve policy. The correct floor depends on the liabilities, flexibility, and exposure actually carried. Do not calculate available capital by subtracting the same cost twice under different labels. Conversely, do not forget that a prepaid program may require substantial future service spending not visible in the current month's bank disbursements.

\textbf{Runway illustration:} ₹7.2 lakh available cash with a separately chosen ₹2.4 lakh reserve floor leaves ₹4.8 lakh discretionary buffer; at ₹80,000 monthly net cash burn this simplistic buffer lasts 6 months. If the reserve floor is already reflected in the cash number, do not subtract it twice. Use a dated cash forecast to locate the actual trough.

\subsection*{Appendix F: Founder Compensation and Conflict-of-Interest Policy}

Write a simple owner-venture agreement before the 1st meaningful distributions or outside investment. Define which work the founder provides, whether salary will be paid, reimbursement rules, how personal and business expenses are separated, what constitutes a capital contribution or loan, and the authority for any distribution. Include any employment-contract and intellectual-property restrictions that must be respected.

A good monthly review shows \textbf{cash compensation drawn}, \textbf{documented unpaid founder-hours}, \textbf{normalized cost at an explicit assumption}, and \textbf{the difference between current cash contribution and normalized economics}. No founder should automatically be treated as entitled to whatever cash remains after enrollment. Nor should a socially ambitious founder be expected to subsidize a perpetually unviable business indefinitely.

If the founder keeps an external job, schedule actual sustainable availability and leave room for ordinary life, recovery, and obligations. A business model that can only be delivered through chronic overwork has a labor-capacity liability even when no accountant has recorded 1.

\subsection*{Appendix G: A Contract-Level Survival Approval Sheet}

Before accepting a material institutional or enterprise deal, complete the following 10-gate decision test:

\begin{table*}[htbp]
\centering
\small
\begin{tabularx}{\textwidth}{lp{7.0cm}X}
\toprule
\textbf{Gate} & \textbf{Evidence Required} & \textbf{Decision Options} \\
\midrule
1. Real buyer & Authorized contracting entity, budget holder, and purchase authority & Accept / clarify / decline \\
2. Defined delivery & Concrete scope, acceptance milestones, standards, and explicit exclusions & Accept / clarify / decline \\
3. Collection route & Verified invoice contact, billing prerequisites, and payment terms & Accept / negotiate / decline \\
4. Direct economics & Realized fee less direct contract-specific costs and mentor labor & Accept / redesign / decline \\
5. Founder capacity & Necessary specialist hours realistically available without burnout & Accept / stage / decline \\
6. Working-capital fit & 1st cash trough modeled under a 60-to-90-day collection delay & Accept / finance / decline \\
7. Customer obligations & Reasonable refund, support, warranty, and remediation terms & Accept / revise / decline \\
8. Shared risk & Absence of operational or IP conflict with existing commitments & Accept / mitigate / decline \\
9. Strategic value & Reusable teaching assets, referenceability, or curriculum accumulated & Accept / test / decline \\
10. Downside limit & Minimum reserve floor remains intact under worst-case client default & Accept / defer / decline \\
\bottomrule
\end{tabularx}
\caption{The 10-gate contract-level survival approval sheet for institutional deals.}
\end{table*}

A contract that fails 1 essential gate should not be rescued by scoring highly on unrelated prestige or brand value. If a strategic experiment is intentionally loss-making, document its maximum loss and the learning asset it is expected to create. That transforms an unpriced risk into a bounded investment thesis.

\subsection*{Appendix H: The 30-Day Financial Survival Sprint}

\begin{enumerate}
\item \textbf{Days 1 to 5: Establish truth.} List every relevant bank account, opening balance, outstanding receivable, payable, refund, loan, and tax obligation. Reconcile bank movements. Identify personal expenses or liabilities that have been mixed into the venture, and document the appropriate resolution with professional advice. \emph{Output: a reliable starting balance sheet and commitments ledger.}

\item \textbf{Days 6 to 10: Map timing.} Build a weekly cash forecast covering 13 weeks. Enter receipts according to probable cash dates, not revenue recognition dates. Enter payroll, supplier, loan, tax, and operating commitments. \emph{Output: 1st projected cash trough, overdue collection list, and minimum reserve policy hypothesis.}

\item \textbf{Days 11 to 15: Make economics truthful.} Compute contribution by product or contract, including acquisition, variable delivery, fixed cohort costs, refunds, and capacity steps. Add a normalized founder-labor analysis separate from current cash expense. \emph{Output: an honest view of what the company earns when it repeats its work.}

\item \textbf{Days 16 to 20: Stress the system.} Delay the largest uncertain collection, reduce demand, add a realistic delivery overrun, and inspect the minimum forecast cash under each scenario. Write preventive actions and spending triggers. \emph{Output: an adverse-case response plan that does not depend on violating customer promises.}

\item \textbf{Days 21 to 25: Put controls in place.} Set payment authorization limits, vendor validation, bank reconciliation, invoice follow-up, tax calendar, and record retention. If staffing is tiny, design proportional compensating reviews rather than pretending full segregation of duties is possible. \emph{Output: reliable basic financial controls.}

\item \textbf{Days 26 to 30: Decide.} Review which projects create contribution, which consume capital, which risks could destroy the business, and which next experiments are both strategically useful and financially affordable. Document the capital-allocation rules and the next 30 days' commercial commitments. \emph{Output: a 1-page survival memo and a calendar for weekly cash reviews and monthly closes.}
\end{enumerate}

The sprint is not a substitute for a qualified accountant or a guarantee of solvency. Its purpose is to move the founder from impressionistic cash management to explicit obligations, dated forecasts, and decisions that can withstand a disappointed customer or a delayed payment.
